% source_equivalent_api_adapters.tex —— engine 实现转写：求解器门面与适配器调度
% 本文件为实现转写章，遵循 docs/modules/README.md 的数学模型规范。
% 所有规则只转写当前源码实际执行的逻辑；锚点禁止行号。

\section{求解器门面、适配器注册与分派}\label{sec:impl-api-adapters}

\subsection{转写范围与口径}\label{subsec:api-scope}

本章转写引擎对外 API 与适配器层的当前实现，核验日期 2026-09-13。覆盖文件：
\begin{itemize}
  \item 门面 \srcpath{src/engine/api/solver.cpp}（\texttt{SolverEngine} 全部成员）；
  \item 适配器基类与注册表
        \srcpath{src/engine/solver/solver\_adapter.cpp}、
        \srcpath{src/engine/solver/adapter\_registry.cpp} 及对应头文件
        \srcpath{include/mipsolvers/engine/solver/solver\_adapter.hpp}、
        \srcpath{include/mipsolvers/engine/api/result.hpp}、
        \srcpath{include/mipsolvers/engine/api/options.hpp}、
        \srcpath{include/mipsolvers/engine/solve\_context.hpp}；
  \item 分派器 \srcpath{src/engine/strategy/dispatcher.cpp}；
  \item 原生适配器
        \srcpath{src/engine/solver/native/native\_adapters.cpp}（LP 选择器
        \texttt{native\_lp\_selector.cpp}、PDLP \texttt{lp/pdlp\_solver.cpp}、IPM 系
        适配器仅述其注册名，实现转写见本手册 LP/IPM 章节）；
  \item 外部适配器 \srcpath{src/engine/solver/external/adapters.cpp}
        （HiGHS/Ipopt/SCIP/Gurobi/CPLEX）；
  \item 策略辅助件 \srcpath{src/engine/strategy/presolve\_manager.cpp}、
        \srcpath{src/engine/strategy/postsolve\_manager.cpp}、
        \srcpath{src/engine/strategy/warmstart.cpp}；
  \item 输入校验 \srcpath{src/engine/util/problem\_validation.cpp}。
\end{itemize}

\begin{boundarynote}
口径约定：（i）本章只描述``数据如何流动与何时回退''，各求解器内部的算法转写
属于各自章节（MILP B\&C 见 \ref{sec:impl-milpbc} 章）。（ii）与用户手册
\texttt{docs/manual/05-solvers-engines.md} \S 5.4 的两处行文差异在
\ref{subsec:api-docdiff} 集中声明；以代码为准。（iii）匿名命名空间内的自由函数
以 \texttt{文件名:函数名} 锚定（不出现 \texttt{::} 前缀），与仓库锚点约定一致。
\end{boundarynote}

\subsection{门面 \texttt{SolverEngine}：规范化、校验与结果映射}\label{subsec:api-facade}

公共求解入口的事务链固定为四步：
\begin{enumerate}
  \item \textbf{矩阵规范化}：对按值传入的 \texttt{ProblemVariant} 原地补齐零行
        矩阵的列数——\texttt{A}/\texttt{Aeq} 行数为 0 但列数不等于 $n$ 时
        resize 成 $0\times n$；QP 的空 Hessian resize 成 $n\times n$；MILP 作用于
        \texttt{linear\_part}。原因是 Eigen 稀疏矩阵默认构造为 $0\times0$，逐列
        迭代的适配器在 \texttt{cols()<n} 时会越界。
        对应源码为
        \srcpath{src/engine/api/solver.cpp:normalize\_lp\_matrices}、
        \srcpath{src/engine/api/solver.cpp:normalize\_problem}；
  \item \textbf{校验}：对规范化后的模型调用 \texttt{validate}（按 variant 分派到
        各模型重载），任一错误即以 \texttt{std::invalid\_argument} 抛出，
        消息为 \texttt{problem\_class\_name} 加分号连接的全部错误。
        对应源码为 \srcpath{src/engine/api/solver.cpp:throw\_if\_invalid}、
        \srcpath{src/engine/util/problem\_validation.cpp:validate}；
  \item \textbf{分派}：构造 \texttt{SolveContext} 并调用
        \texttt{StrategyDispatcher::solve}（见 \ref{subsec:api-dispatch}）；
  \item \textbf{结果映射}：\texttt{SolveResult}$\to$\texttt{api::Result} 为逐字段
        拷贝/移动，无任何重算；\texttt{x}、三组对偶向量与全部统计字段原样透传。
        对应源码为 \srcpath{src/engine/api/solver.cpp:to\_api\_result}、
        \srcpath{src/engine/api/solver.cpp:SolverEngine::solve\_normalized}。
\end{enumerate}
类型化便捷入口（\texttt{solve\_lp}/\texttt{solve\_milp}/\texttt{solve\_qp}/%
\texttt{solve\_conic} 等）只做对应的矩阵规范化后转入统一 \texttt{solve}。
对应源码为 \srcpath{src/engine/api/solver.cpp:SolverEngine::solve\_lp}、
\srcpath{src/engine/api/solver.cpp:SolverEngine::solve}。

输入校验的实现规则（\texttt{problem\_validation.cpp}）：通用变量元数据要求
界有限且 $l_j\le u_j+10^{-12}$（\texttt{kBoundEps}，硬编码）；binary 变量界超出
$[0,1]$ 只产生 warning。LP 要求 $c$ 非空、$\texttt{A.cols}=n$、
\texttt{row\_lhs} 为空或长度等于 \texttt{A.rows} 且每个有限两端行满足
$\mathrm{lhs}_i\le b_i+10^{-12}$；MILP/MINLP 额外要求整数/二进制指标集在界内、
无重复且两集合不交（排序后相邻判等）；Conic 要求
$\texttt{G.rows}=\texttt{h.size}=\texttt{dims.total()}$ 且锥维数非负。
对应源码为
\srcpath{src/engine/util/problem\_validation.cpp:validate\_variable\_meta}、
\srcpath{src/engine/util/problem\_validation.cpp:validate}。

\subsection{资源契约 \texttt{SolveContext}}\label{subsec:api-context}

\texttt{SolveContext} 在构造时固化一次调用级资源合约：时钟从构造点开始；
\texttt{time\_limit\_sec}$>0$ 时记录绝对 deadline。构造期即拒绝非法输入
（时限非有限或为负、线程数为负均抛 \texttt{std::invalid\_argument}）。
\begin{equation}
  \texttt{remaining\_time\_sec}=
  \max\bigl(0,\ \mathrm{deadline}-\mathrm{now}()\bigr),\qquad
  \texttt{stop\_requested}=
  \texttt{stop\_token.stop\_requested()}\ \lor\ \texttt{deadline\_expired()}.
\end{equation}
设计意图（头文件注释）：绝对 deadline 使 $N$ 次回退共享同一份时限，而不是各
拿一份相对时限。对应源码为
\srcpath{include/mipsolvers/engine/solve\_context.hpp:SolveContext}；
其 R1 条款出处为
\srcpath{docs/archive/general\_solver\_performance\_program\_2026-09-13.md}。

\subsection{适配器接口与注册表}\label{subsec:api-registry}

\texttt{SolverAdapter} 对每个问题类提供两个虚函数层：无 context 的
\texttt{solve\_*} 与带 \texttt{SolveContext} 的重载；基类默认实现全部返回
\texttt{unsupported\_result}（\texttt{success=false}、状态为
``Unsupported problem class for adapter <name>: <CLS>''），带 context 的默认实现
委托给无 context 版本（即不感知 deadline 的适配器仍能编译入注册表）。
对应源码为
\srcpath{src/engine/solver/solver\_adapter.cpp:SolverAdapter::unsupported\_result}、
\srcpath{src/engine/solver/solver\_adapter.cpp:SolverAdapter::solve\_lp}。

\texttt{AdapterRegistry} 是一个保序向量：\texttt{register\_adapter} 只追加
（空指针被静默忽略——显式的宽容路径）；\texttt{adapters\_for(cls)} 按注册序返回
全部 \texttt{supports(cls)} 为真的适配器；\texttt{find\_by\_name} 返回首个同名
适配器。对应源码为
\srcpath{src/engine/solver/adapter\_registry.cpp:AdapterRegistry::register\_adapter}、
\srcpath{src/engine/solver/adapter\_registry.cpp:AdapterRegistry::adapters\_for}、
\srcpath{src/engine/solver/adapter\_registry.cpp:AdapterRegistry::find\_by\_name}。

\paragraph{默认注册顺序。}
\texttt{register\_default\_adapters} 通过 \texttt{register\_if\_missing}（按名
幂等，重名跳过）依次注册：

\begin{center}
\small
\begin{tabular}{ll}
\toprule
顺序 & 适配器（注册名）\\
\midrule
1--12 & \texttt{NativeLinear}, \texttt{NativeNewton}, \texttt{NativeAutoLP},
        \texttt{NativeDualSimplex}, \texttt{NativeIPMLP}, \texttt{NativePDLP},
        \texttt{NativeLCQP}, \texttt{NativeIPM}, \texttt{NativeNLP},
        \texttt{NativeConicIPM}, \texttt{StrictHiGHS}, \texttt{NativeBranchAndCut}\\
13--17 & \texttt{Gurobi}, \texttt{CPLEX}, \texttt{HiGHS}, \texttt{Ipopt},
        \texttt{SCIP}\quad（仅当各自 \texttt{available()} 为真）\\
\bottomrule
\end{tabular}
\end{center}

原生 12 个无条件注册；外部 5 个先构造再探测 \texttt{available()}，不可用则不
入表。该顺序即注册表内序，影响 \texttt{adapters\_for} 的尾部枚举序（见
\ref{subsec:api-dispatch}）。注册名均不带 \texttt{Adapter} 后缀。
对应源码为
\srcpath{src/engine/api/solver.cpp:SolverEngine::register\_default\_adapters}、
\srcpath{src/engine/solver/native/native\_lp\_selector.cpp:NativeDualSimplexLPAdapter}。

\texttt{available()} 的探测语义分两类：编译期内嵌库宏
（\texttt{HACDCPF\_HAVE\_HIGHS\_LIB}/\texttt{HACDCPF\_HAVE\_IPOPT}/%
\texttt{HACDCPF\_HAVE\_SCIP\_LIB}）为真时恒可；否则 HiGHS/Ipopt/SCIP 退化为
``可执行文件路径非空''；Gurobi/CPLEX 无子进程路径，仅在对应编译宏且运行时
环境/许可初始化成功（\texttt{env\_ != nullptr}）时可用。注意
\texttt{GurobiAdapter::supports} 把 \texttt{available()} 并入判定，与其他适配器
``支持关系与可用性分离''的写法不同——属实现现状。
对应源码为
\srcpath{src/engine/solver/external/adapters.cpp:HighsAdapter::available}、
\srcpath{src/engine/solver/external/adapters.cpp:IpoptAdapter::available}、
\srcpath{src/engine/solver/external/adapters.cpp:ScipAdapter::available}、
\srcpath{src/engine/solver/external/adapters.cpp:GurobiAdapter::available}、
\srcpath{src/engine/solver/external/adapters.cpp:CplexAdapter::available}、
\srcpath{src/engine/solver/external/adapters.cpp:GurobiAdapter::supports}。

\subsection{候选排序与回退：\texttt{StrategyDispatcher}}\label{subsec:api-dispatch}

\paragraph{默认优先级表。}
每个问题类的内建候选名序列（\texttt{default\_priority\_for}）：
\begin{center}
\small
\begin{tabular}{ll}
\toprule
问题类 & 默认候选序（名前先）\\
\midrule
LE & \texttt{NativeLinear}\\
NLE & \texttt{NativeNewton}\\
LP & \texttt{NativeAutoLP}, \texttt{NativeIPMLP}, \texttt{NativePDLP},
     \texttt{NativeLCQP}, \texttt{HiGHS}\\
QP & \texttt{NativeLCQP}\\
NLP & \texttt{Ipopt}, \texttt{NativeIPM}, \texttt{NativeNLP}\\
MILP & \texttt{StrictHiGHS}, \texttt{HiGHS}, \texttt{NativeBranchAndCut}\\
MINLP & \texttt{Scip}, \texttt{NativeBranchAndCut}\\
CONIC & \texttt{NativeConicIPM}\\
\bottomrule
\end{tabular}
\end{center}
对应源码为
\srcpath{src/engine/strategy/dispatcher.cpp:default\_priority\_for}。

\paragraph{策略排序。}
\texttt{StrategyPolicy::ExternalFirst} 把名字属于
$\{\texttt{HiGHS},\texttt{Ipopt},\texttt{Scip},\texttt{Gurobi},\texttt{CPLEX}\}$
者排前（rank 0）、\texttt{StrictHiGHS} 或以 \texttt{Native} 前缀开头者次之
（rank 1），其余殿后（rank 2）；\texttt{NativeFirst} 对称交换；\texttt{Auto}
不动序。排序为 \texttt{stable\_sort}。按类覆盖
（\texttt{class\_strategy\_policy}）优先于全局策略。
对应源码为
\srcpath{src/engine/strategy/dispatcher.cpp:apply\_policy\_ordering}、
\srcpath{src/engine/strategy/dispatcher.cpp:effective\_policy}。

\paragraph{候选拼装。}
\texttt{candidate\_adapters} 依次并入（\texttt{seen} 集合去重）：
（1）\texttt{preferred\_solver}；（2）按类偏好
（\texttt{set\_solver\_preference} 登记）；（3）策略排序后的默认表；
（4）注册表中其余支持该类的适配器（注册序）。每一步都要求按名找到且
\texttt{supports(cls)} 为真。
对应源码为
\srcpath{src/engine/strategy/dispatcher.cpp:StrategyDispatcher::candidate\_adapters}、
\srcpath{src/engine/strategy/dispatcher.cpp:StrategyDispatcher::set\_solver\_preference}。

\paragraph{回退事务。}
\texttt{StrategyDispatcher::solve} 对候选依次：每次尝试\textbf{之前}检查绝对
deadline 与 stop token（命中则分别返回 \texttt{"Time limit reached before
fallback"}/\texttt{"Cancelled before fallback"}，\texttt{success=false}）；
随后经 \texttt{solve\_with\_adapter} 按 variant 取载荷并调用对应虚函数（载荷
类型不符返回 \texttt{"Problem payload mismatch ..."}）；结果
\texttt{stats.success=true} 即返回；否则记录为最近失败，且
\texttt{allow\_fallback=false} 时立即返回该失败。全部候选失败时返回
\textbf{最后一次}失败的结果（而非第一个）。候选为空返回
\texttt{"No adapter registered for class ..."}。
对应源码为
\srcpath{src/engine/strategy/dispatcher.cpp:StrategyDispatcher::solve}、
\srcpath{src/engine/strategy/dispatcher.cpp:solve\_with\_adapter}。

\begin{boundarynote}
回退的两个实现现状，引用时不得读成更强合约：
（i）回退只在 \texttt{success=false} 上触发；\texttt{success=true} 但
\texttt{status} 表明次优/截断的结果不会触发下一个候选。（ii）注册表尾部枚举
会把\textbf{所有}注册且支持该类的适配器追加为末位候选，包括 Gurobi/CPLEX
（只要 \texttt{available()} 使它们入了注册表）；它们只是不在
\texttt{default\_priority\_for} 的内建优先级表内。用户手册 \S 5.4.2
``Gurobi 与 CPLEX 不参与普通自动回退''一句与该尾部路径不符，见
\ref{subsec:api-docdiff}。
对应源码为
\srcpath{src/engine/strategy/dispatcher.cpp:StrategyDispatcher::solve}、
\srcpath{src/engine/strategy/dispatcher.cpp:StrategyDispatcher::candidate\_adapters}。
\end{boundarynote}

\subsection{原生适配器的适配层}\label{subsec:api-native}

\paragraph{LE/NLE/NLP 轻量适配器。}
\texttt{NativeLinear}（LE）直接稀疏 LU：分析-分解-求解后以原系统残差
$\|Ax-b\|_\infty$ 复核，状态 \texttt{"Optimal"}。
\texttt{NativeNewton}（NLE）是阻尼 Newton：正则化阶梯
$\lambda: \texttt{regularization0}\to 10^{-10}\to \times 10$（至多 4 次尝试，
硬编码）解 $(J+\lambda I)\,\Delta x=-f$，回溯线搜索要求
$\|f(x+\alpha\Delta x)\|_\infty\le\|f(x)\|_\infty$，步长回退因子
\texttt{step\_backoff}。
\texttt{NativeNLP} 是固定罚参数的二次罚 Newton（非 IPM）：罚梯度
$\nabla\phi=\nabla f+\rho J_g^{T}g+\rho J_h^{T}h^{+}$，罚 Hessian
$H=\nabla^2 f+\rho J_g^{T}J_g+\rho (J_h^{\mathrm{act}})^{T}J_h^{\mathrm{act}}$
再加对角正则 \texttt{regularization0}；收敛要求
$\|\nabla\phi\|_\infty$、$\|g\|_\infty$、$\|h^{+}\|_\infty$ 同时
$\le\texttt{tol\_grad}$，或 $\|\Delta x\|_\infty\le\texttt{tol\_step}$；
接受判据为罚函数不增。头文件级注释明确它不是 \texttt{NativeIPM}。
对应源码为
\srcpath{src/engine/solver/native/native\_adapters.cpp:NativeLinearAdapter::solve\_le}、
\srcpath{src/engine/solver/native/native\_adapters.cpp:NativeNewtonAdapter::solve\_nle}、
\srcpath{src/engine/solver/native/native\_adapters.cpp:NativeNLPAdapter::solve\_nlp}。

\paragraph{B\&C 适配器与 LP 快路径。}
\texttt{NativeBranchAndCut} 支持 MILP 与 MINLP。MILP 入口先走 LP 快路径：当
\texttt{binary\_idx} 与 \texttt{integer\_idx} 皆空时直接以单纯形求解
（迭代上限 $\max(\texttt{max\_lp\_iter},20000)$，容差
$\max(10^{-8},0.1\,\texttt{lp\_tol})$），并从最优基提取约束对偶：标准形空间
\begin{equation}
  y = B^{-T}c_B,\qquad
  d_i = s_{\mathrm{sense}}\cdot\mathrm{row\_sign}_i\cdot y_i,\quad
  s_{\mathrm{sense}}=\begin{cases}-1,&\text{Minimize}\\+1,&\text{Maximize}\end{cases}
\end{equation}
（标准形内部为最大化，故最小化要再翻一次符号）；优先用稠密基逆，缺省时用
稀疏 BTRAN，再缺省则对偶置零——显式回退链。布局统一为
$[\,\text{不等式行对偶}\;|\;\text{等式行对偶}\,]$。
带整数的模型转入 \texttt{solve\_milp\_bc}（见 \ref{sec:impl-milpbc} 章）。
context 版入口把绝对 deadline 的剩余时间与线程预算写进 \texttt{BCOptions}
副本（\texttt{time\_limit\_sec}/\texttt{num\_threads}）。
对应源码为
\srcpath{src/engine/solver/native/native\_adapters.cpp:NativeBranchAndCutAdapter::solve\_milp}、
\srcpath{src/engine/solver/native/milp/bc/api.cpp:solve\_milp\_bc}。

\paragraph{StrictHiGHS 选项改造。}
\texttt{StrictHiGHS} 是同名 B\&C 内核的选项变体（委托给
\texttt{NativeBranchAndCutAdapter} 执行）：\texttt{make\_strict\_highs\_production\_%
options} 固定 \texttt{lp\_kernel\_backend=HiGHS}、\texttt{strict\_highs\_mip\_%
contract=true}、\texttt{auto\_highs\_root\_pipeline=true}，开启
\texttt{highs\_mip\_run\_zi\_round}/\texttt{highs\_mip\_run\_shifting}，强制
presolve on 并把 \texttt{highs\_presolve\_substitution\_maxfillin} 提到 30，
\texttt{highs\_mip\_lp\_age\_limit=30}、\texttt{highs\_mip\_detect\_symmetry=true}。
\texttt{make\_strict\_highs\_problem\_options} 再按模型规模决定是否把根 LP
切成 IPM：当列数/行数进入
$[\texttt{min\_cols},\texttt{max\_cols})\times[\texttt{min\_rows},\texttt{max\_rows})$
窗口、全局开关
\texttt{highs\_strict\_auto\_ipm\_root\_for\_large\_models} 开启、且时限满足
最小预算
\begin{equation}
  T_{\min}=\max\Bigl(\texttt{min\_time\_sec},\
  \frac{n_{\mathrm{cols}}\cdot\texttt{secs\_per\_kcol}}{1000}\Bigr)
\end{equation}
（默认 \texttt{secs\_per\_kcol}$=2.0$ 时 6360 列对应 $30$\,s、26400 列对应
$52.8$\,s——代码注释中的两个参考点）时，置 \texttt{highs\_mip\_lp\_solver="ipm"}
并把 crossover 默认开。显式调用者设置不被该规模门覆盖。
对应源码为
\srcpath{src/engine/solver/native/native\_adapters.cpp:make\_strict\_highs\_production\_options}、
\srcpath{src/engine/solver/native/native\_adapters.cpp:make\_strict\_highs\_problem\_options}、
\srcpath{src/engine/solver/native/native\_adapters.cpp:StrictHighsBranchAndCutAdapter::solve\_milp}。

\subsection{外部适配器的模型映射与结果解析}\label{subsec:api-external}

五个外部适配器共享同一骨架：校验 $\to$ \texttt{available()} 门 $\to$
等价模型映射 $\to$ 内嵌库调用（编译期启用时优先）或子进程回退 $\to$
结果解析 $\to$ 统一 \texttt{SolveResult}。不可用/失败时返回
\texttt{unavailable\_result}（\texttt{success=false}、状态以
\texttt{"Unavailable:"} 等前缀说明原因）。
对应源码为
\srcpath{src/engine/solver/external/adapters.cpp:unavailable\_result}。

\paragraph{HiGHS。}
支持 LP/MILP。MILP 映射：把 \texttt{integer\_idx}/\texttt{binary\_idx} 提升为
\texttt{linear\_part} 上的变量类型，binary 界夹取到 $[0,1]$；
\texttt{initial\_solution} 尺寸匹配时作为 MIP start 传入。内嵌路径经
\texttt{solve\_lp\_with\_embedded\_highs}，调用前先
\texttt{Highs::resetGlobalScheduler(true)} 以消除先前 MILP 求解遗留的全局线程
调度器配置（注释说明原因）。子进程回退：写临时 MPS（列名 \texttt{X1..Xn}，
1-indexed）+ 选项文件（LP 写 \texttt{time\_limit}/\texttt{threads}/%
\texttt{random\_seed}；MILP 另写死 \texttt{mip\_rel\_gap = 1e-4}——硬编码），
调 \texttt{--model\_file/--solution\_file/--options\_file}，解析 .sol 的
\texttt{Model status} 与日志；\texttt{success}$\iff$状态为
\texttt{Optimal} 或 \texttt{Feasible}，解值按 \texttt{X<i+1>} 名回映射，缺失列
补 0。临时文件求解后删除。
对应源码为
\srcpath{src/engine/solver/external/adapters.cpp:HighsAdapter::solve\_lp\_impl}、
\srcpath{src/engine/solver/external/adapters.cpp:HighsAdapter::solve\_milp\_impl}。
另有确定性定价入口 \texttt{solve\_pricing\_lp}：只走内嵌库，参数非法直接抛
\texttt{std::invalid\_argument}。
对应源码为
\srcpath{src/engine/solver/external/adapters.cpp:HighsAdapter::solve\_pricing\_lp}。

\paragraph{SCIP。}
支持 MILP/MINLP。MILP 走 MPS 交换（同样的整数提升与 binary 夹取）；内嵌路径
\texttt{SCIPreadProb} 读入后求解，固定参数 \texttt{limits/gap=1e-3}、
\texttt{limits/time=300.0}（硬编码，注释说明与原生 UC 容差对齐）、
\texttt{display/verblevel=0}；有解即 \texttt{success=true}，
\texttt{status} 依 \texttt{SCIPgetStatus} 区分 \texttt{"Solved"}/%
\texttt{"Feasible (limit)"}/\texttt{"Infeasible"}，gap 取
\texttt{SCIPgetGap}；解值经 \texttt{SCIPgetOrigVars}（原始空间，而非 presolve
后的变换变量——注释说明后者名字对不上 MPS 导出）按名回映射。子进程回退用
\texttt{-c "set limits gap 0.001" -c read -c optimize -c "write solution"} 脚本。
MINLP：无符号目标时走 NLP 松弛 + 舍入回退（整数取整、binary 按 $0.5$ 阈值、
再夹回界内，$\|g\|_\infty$/$h^+$ 复核 $\le 10^{-4}$ 才算成功——硬编码阈值，
且这只是启发式，不是全局证明）；有符号目标时写 PIP 文件交给 SCIP。
对应源码为
\srcpath{src/engine/solver/external/adapters.cpp:ScipAdapter::solve\_milp}、
\srcpath{src/engine/solver/external/adapters.cpp:ScipAdapter::solve\_minlp}。

\paragraph{Ipopt。}
仅 NLP，且\textbf{只有内嵌库路径真正求解}：无内嵌时即使有可执行文件也返回
\texttt{"Unavailable: external Ipopt executable adapter is disabled by default"}
（显式禁用路径）。求解前：\texttt{x0} 尺寸不符时重构为界内中点/端点的默认
点；缺 \texttt{f}/\texttt{grad} 或有 \texttt{g} 无 \texttt{jac\_g}、有
\texttt{h} 无 \texttt{jac\_h} 直接 Unavailable；primal-dual warm start 需维度
匹配、全部有限且界对偶非负，否则报错。选项映射：
\texttt{hessian\_approximation} 按有无 Lagrangian Hessian 回调取
\texttt{exact}/\texttt{limited-memory}；\texttt{bound\_relax\_factor=0.0}
（注释：默认界松弛会在大乘子下破坏原模型互补性）；warm start 统一推送量
默认 $10^{-8}$。容差经 \texttt{positive\_finite\_or} 兜底（非有限或非正回退
默认），并强制 Ipopt 的不变量
$\texttt{acceptable\_tol}\ge\texttt{tol}$（对偶/约束/互补四组同法）。
结果由 \texttt{CallbackTNLP::build\_result} 从 \texttt{OptimizeTNLP} 状态装配。
对应源码为
\srcpath{src/engine/solver/external/adapters.cpp:IpoptAdapter::solve\_nlp}。

\paragraph{CPLEX。}
仅 MILP，Callable Library 内嵌（无子进程路径）。构造期打开环境并固定
\texttt{ScreenOutput=off}、\texttt{TimeLimit}、\texttt{MIPGap}、\texttt{Threads}、
\texttt{RandomSeed} 五个参数，任一失败即记 \texttt{initialization\_error\_} 且
环境置空。模型映射的两处精确转写：（i）双端行展开——每个不等式行按其有限端
拆成至多两行（上端 \texttt{'L'}、下端 \texttt{'G'}），等式行为 \texttt{'E'}，
列主序打包时同一列元素向所属各展开行各写一份；行数或非零元数超
\texttt{int} 上限时失败返回。（ii）列类型 \texttt{'I'}/\texttt{'B'}，binary
界夹取到 $[0,1]$，非有限界映射为 $\pm$\texttt{CPX\_INFBOUND}。结果解析：
\texttt{success} 以是否取到解为准（\texttt{CPXgetx} 失败会把 success 翻回
false），状态字符串原样取 \texttt{CPXgetstatstring}，gap 取
\texttt{CPXgetmiprelgap}；\texttt{cplex\_status\_proven}/\texttt{optimal}/%
\texttt{timed\_out} 三个谓词把 \texttt{CPXMIP\_*} 状态码分类进
\texttt{thread\_local} 诊断 \texttt{cplex\_solve\_info}。
对应源码为
\srcpath{src/engine/solver/external/adapters.cpp:CplexAdapter::solve\_milp}、
\srcpath{src/engine/solver/external/adapters.cpp:CplexAdapter::CplexAdapter}；
双端行展开的正确性论证见
\srcpath{docs/archive/cplex\_callable\_library\_integration\_2026-09-12.md}。

\paragraph{Gurobi。}
支持 LP/QP/MILP（\texttt{supports} 依赖 \texttt{available()}，见
\ref{subsec:api-registry}）。LP 映射：最大化经目标取负统一为
\texttt{GRB\_MINIMIZE}（$\texttt{obj\_sign}=-1$）；不等式行按下标
$|a_{ij}|>10^{-15}$ 过滤后逐行 \texttt{GRBaddconstr}（双端行拆成
\texttt{GRB\_LESS\_EQUAL} 与 \texttt{GRB\_GREATER\_EQUAL} 两条，$10^{-15}$
为硬编码）；行按\textbf{稠密提取}构造（代码保留了稀疏列主序迭代的空循环
与注释，说明为何不用它——已知的历史写法残留，不影响结果）。两个特化入口：
\texttt{solve\_pricing\_lp} 按规模固定方法（$n\ge10^{6}$ 用
\texttt{method=2}/\texttt{threads=8}/\texttt{crossover=0}，否则
\texttt{method=1}/单线程，硬编码阈值），并把 \texttt{OptimalityTol=1e-8} 写进
环境；\texttt{solve\_relaxation\_lp} 是不精确 barrier 入口
（\texttt{BarConvTol} 取参，合法域 $[10^{-8},10^{-2}]$，之外抛异常；注释明确
禁止用它认证价格或原始可行性）。
对应源码为
\srcpath{src/engine/solver/external/adapters.cpp:GurobiAdapter::solve\_lp}、
\srcpath{src/engine/solver/external/adapters.cpp:GurobiAdapter::solve\_pricing\_lp}、
\srcpath{src/engine/solver/external/adapters.cpp:GurobiAdapter::solve\_relaxation\_lp}。

\subsection{策略辅助件：presolve/postsolve/warmstart 管理器}\label{subsec:api-strategy-aux}

\begin{boundarynote}
接线现状（必须在引用处保持可见）：\texttt{PresolveManager}/%
\texttt{PostsolveManager}/\texttt{WarmStartManager} 是\textbf{独立工具类}，
当前未被 \texttt{StrategyDispatcher::solve} 调用——分派器直接调用适配器，主流程
的 presolve 发生在各求解器内部（MILP 的 presolve 见 \ref{sec:impl-milpbc} 章）。
\texttt{PresolveManager::should\_presolve} 恒返回 \texttt{false}，
\texttt{PresolveManager::presolve} 退化为恒等映射 \texttt{native\_presolve}
（只统计行列数，不做约简）；\texttt{PresolveManager::scale} 对
\texttt{"none"}/\texttt{"identity"} 之外的 scaling 直接抛
\texttt{std::logic\_error}（无通用缩放后端）。这些类的完整实现仅在测试中被
消费。对应源码为
\srcpath{src/engine/strategy/presolve\_manager.cpp:PresolveManager::should\_presolve}、
\srcpath{src/engine/strategy/presolve\_manager.cpp:PresolveManager::native\_presolve}、
\srcpath{src/engine/strategy/presolve\_manager.cpp:PresolveManager::scale}、
\srcpath{src/engine/strategy/dispatcher.cpp:StrategyDispatcher::solve}。
\end{boundarynote}

\texttt{PostsolveManager} 实现的三步还原（给定 \texttt{PresolveMapping}）：
\begin{enumerate}
  \item 固定变量还原：按 \texttt{col\_mapping} 把约简解散布回原列，再写入
        \texttt{fixed\_variables}；映射重复、固定列与活跃列重叠、或还原后存在
        非有限分量（未覆盖列）均抛 \texttt{std::invalid\_argument}；
  \item 平移还原：$x_{\mathrm{orig}} \mathrel{+}= \mathrm{shift}$（逐条
        \texttt{variable\_shifts}）；
  \item 反缩放：$x_j \mathrel{*}= s^{\mathrm{col}}_j$、
        $y_i \mathrel{/}= s^{\mathrm{row}}_i$，尺度必须有限且为正。
\end{enumerate}
一致性校验 \texttt{validate} 比较两结果的 \texttt{success}、解维数与有限性，
目标差容差为 $\texttt{feasibility\_tol}\cdot\max\{1,|f_1|,|f_2|\}$。
对应源码为
\srcpath{src/engine/strategy/postsolve\_manager.cpp:PostsolveManager::postsolve}、
\srcpath{src/engine/strategy/postsolve\_manager.cpp:PostsolveManager::restore\_fixed\_variables}、
\srcpath{src/engine/strategy/postsolve\_manager.cpp:PostsolveManager::unscale}、
\srcpath{src/engine/strategy/postsolve\_manager.cpp:PostsolveManager::validate}。

\texttt{WarmStartManager} 的规则：\texttt{is\_compatible} 逐槽检查 primal/约束
对偶/界对偶向量的维数；\texttt{extract\_from\_result} 只在尺寸匹配时搬运各槽；
\texttt{merge} 取各槽首个非空候选，\texttt{reuse\_root\_factorization} 取或；
\texttt{scale} 的坐标规则为 $x'=x/s_x$、$y'=y\odot s_c$、界对偶 $\odot\, s_x$；
\texttt{validate} 拒绝非有限 hint。
对应源码为
\srcpath{src/engine/strategy/warmstart.cpp:WarmStartManager::is\_compatible}、
\srcpath{src/engine/strategy/warmstart.cpp:WarmStartManager::extract\_from\_result}、
\srcpath{src/engine/strategy/warmstart.cpp:WarmStartManager::merge}、
\srcpath{src/engine/strategy/warmstart.cpp:WarmStartManager::scale}。

\subsection{结果契约}\label{subsec:api-result}

内部 \texttt{SolveResult}/\texttt{SolveStats} 与公共 \texttt{api::Result}/%
\texttt{api::Stats} 的字段一一对应（\texttt{to\_api\_result} 不透传的内部字段
仅有 \texttt{primal\_ray} 与稀疏代数/原生单纯形遥测等实现侧计数——未被公共
API 消费的字段在此声明）。契约要点：
\begin{itemize}
  \item \texttt{success} 只表示``该适配器成功终止''：B\&C 在有时限可行
        incumbent 时亦为 true（状态 \texttt{kFeasibleIncumbentFound} 系），
        与``已证最优''无关；验收必须同时看 \texttt{status}、可行度、
        \texttt{mip\_gap}；
  \item \texttt{status} 为人类可读字符串，无枚举约束；同义词集合以
        \texttt{bc\_status} 常量与各适配器字面量为准；
  \item LP 对偶布局跨适配器统一为 $[\,\text{不等式行}\;|\;\text{等式行}\,]$，
        界对偶 \texttt{box\_dual\_lb}/\texttt{box\_dual\_ub} 每列一个；
        MILP 等无证书路径置空；
  \item 不可行证书 \texttt{farkas\_ray}/\texttt{farkas\_ray\_eq}（配合
        \texttt{has\_farkas\_certificate}）：对
        $\min c^Tx\ \mathrm{s.t.}\ Ax\le b,\ A_{\mathrm{eq}}x=b_{\mathrm{eq}},\
        l\le x\le u$ 的不可行 LP，满足 $y\ge0$、$A^Ty+A_{\mathrm{eq}}^Ty_{eq}=0$、
        $b^Ty+b_{\mathrm{eq}}^Ty_{eq}<0$；
  \item \texttt{certified\_dual\_bound} 仅由原生对偶单纯形内核在中断时发布；
        中断结果的 \texttt{objective} 仍是未认证的工作值（头文件注释明确）。
\end{itemize}
对应源码为
\srcpath{include/mipsolvers/engine/solver/solver\_adapter.hpp}、
\srcpath{include/mipsolvers/engine/api/result.hpp:Stats}、
\srcpath{src/engine/api/solver.cpp:to\_api\_result}。

\subsection{与用户手册行文的两处差异}\label{subsec:api-docdiff}

以下差异按``以代码为准''声明，供手册维护侧对照：

\begin{enumerate}
  \item \texttt{docs/manual/05-solvers-engines.md} \S 5.4.2 把调度流程写成
        ``预处理 $\to$ 适配器调用 $\to$ 回退 $\to$ 后处理''六步；当前
        \texttt{StrategyDispatcher::solve} 不含门面级 presolve/postsolve 阶段
        （\texttt{PresolveManager} 未被接线，见
        \ref{subsec:api-strategy-aux}），实际为
        ``候选构造 $\to$ 逐个适配器调用（带绝对 deadline 门）$\to$ 成功即返回
        或回退''。模型规范化与校验在 \texttt{SolverEngine::solve} 一侧完成。
        对应源码为
        \srcpath{src/engine/strategy/dispatcher.cpp:StrategyDispatcher::solve}、
        \srcpath{src/engine/api/solver.cpp:SolverEngine::solve}；
  \item 同节称``Gurobi 与 CPLEX 不参与普通自动回退''；它们确实不在
        \texttt{default\_priority\_for} 内建表内，但一旦注册（可用性门通过），
        \texttt{candidate\_adapters} 的注册表尾部枚举会把它们追加为末位回退
        候选。对应源码为
        \srcpath{src/engine/strategy/dispatcher.cpp:StrategyDispatcher::candidate\_adapters}、
        \srcpath{src/engine/strategy/dispatcher.cpp:default\_priority\_for}。
\end{enumerate}
